% !TEX root = ../../main.tex

\subsection{المرحلة الأولى: النظام بدون ترشيح}
في هذه المرحلة، تم تشغيل الدارة في حالتها الأصلية حيث يغذي المنبع حمل القنطرة الديودية مباشرة عبر ممانعات الشبكة والكبل.

\begin{figure}[H]
  \vspace{15pt}
    \centering
    \includegraphics[width=1.05\linewidth]{00 - Circuit without Filters.png}
    \caption{المخطط الكامل للدارة بدون مرشحات}
    \label{fig:baseline_circuit}
\end{figure}

أظهرت نتائج التحليل في هذه المرحلة تشوهاً كبيراً جداً في موجة التيار، حيث بلغت نسبة \el{\boldmath $THD_I = 24.76\%$}, بينما كانت موجة الجهد أقل تأثراً بنسبة \el{\boldmath $THD_V = 9.55\%$} نتيجة وجود ممانعات الشبكة التي تعمل كمخمدات طبيعية.

\newpage
\begin{figure}[H]
  \vspace{-10pt}
  \centering
  \includegraphics[width=1.05\textwidth]{01 - CurrentData without filters.png}
    \caption{\small نتائج تحليل \el{FFT} للتيار قبل الترشيح - تظهر \el{THD} بنسبة $24.76\%$}
    \label{fig:01}
\end{figure}
\begin{figure}[H]
  \centering
  \includegraphics[width=1.05\textwidth]{02 - VoltageData without filters.png}
    \caption{\small نتائج تحليل \el{FFT} للجهد قبل الترشيح - تظهر \el{THD} بنسبة $9.55\%$}
    \label{fig:02}
\end{figure}

\newpage

